\chapter{Iterative Gegenbeispiel-Konstruktion innerhalb einer SCC}
\label{ch:intra-scc-cegar}

\section{Problemstellung}
\label{sec:intra-scc-problem}

In diesem Kapitel betrachten wir den Spezialfall, dass das gesamte
Argumentationssystem \(F=(A,R)\) aus einer einzigen stark zusammenhängenden
Komponente besteht und die zu testenden Mengen \(X,Y,Z\subseteq A\) paarweise
disjunkt innerhalb dieser SCC liegen.

Statt \(\textit{Ind}_{\mathrm{st}}\) direkt über die in
Kapitel~\ref{ch:stable} angegebene QBF-Kodierung zu entscheiden, verfolgen wir
hier den Ansatz, ein Gegenbeispiel \emph{iterativ} zu konstruieren. Ein
Gegenbeispiel zeigt unmittelbar
\(X\not\bot_{\mathrm{st}} Y \mid Z\); bricht die Konstruktion erfolglos ab,
gilt umgekehrt \(X \bot_{\mathrm{st}} Y \mid Z\).

\begin{Definition}
    \label{def:gegenbsp-st}
    Sei \(F=(A,R)\) ein AF und seien \(X,Y,Z\subseteq A\) paarweise
    disjunkt. Ein \emph{Gegenbeispiel} zur \textit{stable}-Unabhängigkeit
    \(X \bot_{\mathrm{st}} Y \mid Z\) ist ein Paar
    \(\lambda_1,\lambda_2 \in \varLambda_{\mathrm{st}}(F)\) mit
    \begin{enumerate}
        \item \(\lambda_1|_Z = \lambda_2|_Z\),
        \item es existiert kein \(\lambda_3\in\varLambda_{\mathrm{st}}(F)\)
              mit \(\lambda_3|_X = \lambda_1|_X\) und
              \(\lambda_3|_{Y\cup Z} = \lambda_2|_{Y\cup Z}\).
    \end{enumerate}
\end{Definition}

\section{Triviale Fälle: SCC ist ein einzelner Zykel}
\label{sec:intra-scc-zykel}

Bevor wir den allgemeinen Algorithmus formulieren, halten wir zwei
Sonderfälle fest, die in konstanter bzw. linearer Zeit entschieden werden
können.

Sei \(F=(A,R)\) ein AF mit \(|A|=n\) und \(R\) bestehe genau aus dem Zykel
\(a_0\to a_1 \to \dots \to a_{n-1} \to a_0\).

\begin{Lemma}
    \label{lem:zykel-st}
    \begin{enumerate}
        \item Ist \(n\) ungerade, so gilt
              \(\varLambda_{\mathrm{st}}(F)=\emptyset\). Damit ist
              \(X \bot_{\mathrm{st}} Y \mid Z\) für jedes Tripel
              \((X,Y,Z)\) vakuum erfüllt.
        \item Ist \(n\) gerade, so existieren genau zwei
              \textit{stable}-Labellings: die alternierende Färbung
              \(\lambda^*\) mit \(\lambda^*(a_i)=\mathit{in}\) für
              gerades \(i\), und ihre Komplementärfärbung
              \(\overline{\lambda^*}\). Ein Gegenbeispiel existiert
              genau dann, wenn
              \(\lambda^*|_Z = \overline{\lambda^*}|_Z\) und
              \(\lambda^*|_Y \neq \overline{\lambda^*}|_Y\).
    \end{enumerate}
\end{Lemma}

Der gerade Fall verlangt insbesondere, dass entweder \(Z=\emptyset\) ist oder
\(Z\) nur Argumentpaare \(a_i,a_j\) mit \(i+j\) ungerade enthält. Sobald ein
einziges \(z\in Z\) festgelegt wird, fixiert es bereits die gesamte
alternierende Färbung.

\section{Iterativer Algorithmus}
\label{sec:intra-scc-iter}

Der Algorithmus arbeitet ausschließlich auf der aussagenlogischen
Kodierung \(\mathrm{st}_F(I)\) aus Definition~\ref{def:st-pl}. Statt
zwei vollständige Labellings konkret aufzubauen, ratet der Algorithmus
nur die Belegungen auf den ausgezeichneten Mengen \(X\), \(Y\), \(Z\)
und überlässt jede Realisierbarkeitsfrage einem SAT-Aufruf auf
\(\mathrm{st}_F(I)\) mit fixierten Assumptions. Es genügt zu wissen,
\emph{dass} eine Belegung auf \(X\cup Y\cup Z\) zu mindestens einem
stable-Labelling fortgesetzt werden kann --- die konkrete
Vervollständigung außerhalb wird nicht benötigt.

\subsection{Realisierbarkeitsabfrage}
\label{subsec:intra-scc-realize}

Sei \(F=(A,R)\) ein AF, \(W\subseteq A\) eine Teilmenge und
\(\pi : W \to \{\mathit{in},\mathit{out}\}\) eine Belegung von \(W\).
Wir schreiben
\[
    \mathrm{realisierbar}(\pi)
    \quad:\Longleftrightarrow\quad
    \exists \lambda \in \varLambda_{\mathrm{st}}(F):
    \lambda|_W = \pi.
\]
Diese Abfrage ist ein einzelner SAT-Aufruf auf \(\mathrm{st}_F(I)\) mit
den Assumptions
\(\{\, i_a \mid \pi(a) = \mathit{in}\,\} \cup
\{\, \neg i_a \mid \pi(a) = \mathit{out}\,\}\).
Die zurückgelieferte vollständige Belegung wird im Algorithmus nicht
weiter benötigt; einzig die SAT/UNSAT-Antwort zählt.

\subsection{Drei Abfragen pro Kandidat}
\label{subsec:intra-scc-three-queries}

Ein Gegenbeispiel ist nach Definition~\ref{def:gegenbsp-st} festgelegt
durch eine Z-Belegung \(\zeta\) und zwei Restriktionen
\(\pi^1_{X\cup Y}, \pi^2_{X\cup Y}\) von \(\lambda_1, \lambda_2\) auf
\(X\cup Y\) mit
\begin{itemize}
    \item \(\pi^1_X \neq \pi^2_X\) und \(\pi^1_Y \neq \pi^2_Y\),
    \item \(\mathrm{realisierbar}(\zeta \cup \pi^i_{X\cup Y})\) für
          \(i=1,2\),
    \item \(\lnot\mathrm{realisierbar}(\zeta \cup \pi^1_X \cup \pi^2_Y)\).
\end{itemize}
Die ersten beiden Punkte stellen sicher, dass \(\lambda_1, \lambda_2\)
als stable-Labellings überhaupt existieren; der dritte Punkt ist die
eigentliche Gegenbeispiel-Bedingung (die ``Brücken''-Anfrage nach
\(\lambda_3\) ist unerfüllbar). Pro Kandidat fallen damit genau drei
SAT-Aufrufe an: zwei SAT-Tests für die Existenz von \(\lambda_1,
\lambda_2\), ein UNSAT-Test für die Nicht-Existenz von \(\lambda_3\).

\subsection{Algorithmus}
\label{subsec:intra-scc-algo}

\begin{enumerate}
    \item \label{step:enum-z}
        \textbf{Enumeration über \(Z\).} Für jede Belegung
        \(\zeta : Z \to \{\mathit{in},\mathit{out}\}\):

    \begin{enumerate}
        \item \label{step:check-z-real}
            \textbf{Z-Realisierbarkeit.} Prüfe
            \(\mathrm{realisierbar}(\zeta)\). Ist \(\zeta\) nicht
            realisierbar, fahre mit dem nächsten \(\zeta\) fort.

        \item \label{step:enum-xy}
            \textbf{Enumeration über \(X\cup Y\).} Für jedes geordnete
            Paar
            \((\pi^1_{X\cup Y}, \pi^2_{X\cup Y})\) mit
            \(\pi^1_X \neq \pi^2_X\) und \(\pi^1_Y \neq \pi^2_Y\):

        \begin{enumerate}
            \item \label{step:lambda1}
                Prüfe \(\mathrm{realisierbar}(\zeta \cup
                \pi^1_{X\cup Y})\). Falls UNSAT, weiter zum nächsten
                Paar.

            \item \label{step:lambda2}
                Prüfe \(\mathrm{realisierbar}(\zeta \cup
                \pi^2_{X\cup Y})\). Falls UNSAT, weiter zum nächsten
                Paar.

            \item \label{step:bridge}
                Prüfe \(\mathrm{realisierbar}(\zeta \cup \pi^1_X \cup
                \pi^2_Y)\) (\,die Brücken-Anfrage\,). Falls SAT, weiter
                zum nächsten Paar. Falls UNSAT, gib
                \((\zeta, \pi^1_{X\cup Y}, \pi^2_{X\cup Y})\) als
                Gegenbeispiel zurück und terminiere mit
                \(X \not\bot_{\mathrm{st}} Y \mid Z\).
        \end{enumerate}
    \end{enumerate}

    \item \label{step:exhaust}
        Wurden alle Tripel \((\zeta, \pi^1_{X\cup Y}, \pi^2_{X\cup Y})\)
        ohne Erfolg durchlaufen, so gilt \(X \bot_{\mathrm{st}} Y \mid Z\).
\end{enumerate}

Der gesamte Algorithmus besteht aus Schleifen über kombinatorische
Belegungen und Aufrufen einer einzigen Hilfsroutine
\(\mathrm{realisierbar}(\cdot)\). Es werden weder partielle Labellings
verwaltet noch Vervollständigungen außerhalb \(X\cup Y\cup Z\)
explizit konstruiert. Die strukturelle Information, die in Schritt 8
(``ein Gegenbeispiel wurde gefunden'') geliefert wird, sind die drei
Restriktionen \(\zeta\), \(\pi^1_{X\cup Y}\), \(\pi^2_{X\cup Y}\). Aus
ihnen lässt sich bei Bedarf nachträglich ein konkretes Paar
\((\lambda_1, \lambda_2)\) durch zwei weitere SAT-Aufrufe gewinnen
(die Vervollständigungen waren ja in den Schritten~\ref{step:lambda1}
und~\ref{step:lambda2} schon einmal als SAT-Modelle verfügbar; eine
Implementierung kann sie zwischenspeichern, wenn sie das Zeugnis
braucht).

\subsection{Korrektheit}
\label{subsec:intra-scc-correctness}

\begin{Lemma}
    \label{lem:algo-korrekt}
    Der Algorithmus aus \ref{subsec:intra-scc-algo} entscheidet
    \(\textit{Ind}_{\mathrm{st}}\) korrekt.
\end{Lemma}

\begin{proof}
    \emph{Soundness.} Gibt der Algorithmus in Schritt~\ref{step:bridge}
    ein Gegenbeispiel zurück, so sind nach den Schritten~\ref{step:lambda1}
    und~\ref{step:lambda2} \(\lambda_1, \lambda_2 \in
    \varLambda_{\mathrm{st}}(F)\) mit \(\lambda_i|_{X\cup Y} =
    \pi^i_{X\cup Y}\) und \(\lambda_i|_Z = \zeta\) gegeben. Schritt~7
    sagt UNSAT, also existiert kein \(\lambda_3 \in
    \varLambda_{\mathrm{st}}(F)\) mit \(\lambda_3|_X = \pi^1_X =
    \lambda_1|_X\) und \(\lambda_3|_{Y\cup Z} = \pi^2_Y \cup \zeta =
    \lambda_2|_{Y\cup Z}\). Damit ist \((\lambda_1, \lambda_2)\) ein
    Gegenbeispiel im Sinne von Definition~\ref{def:gegenbsp-st}.

    \emph{Completeness.} Angenommen \(X \not\bot_{\mathrm{st}} Y \mid Z\).
    Dann existiert ein Gegenbeispiel \((\lambda_1, \lambda_2)\) nach
    Definition~\ref{def:gegenbsp-st}. Sei
    \(\zeta := \lambda_1|_Z = \lambda_2|_Z\),
    \(\pi^i_{X\cup Y} := \lambda_i|_{X\cup Y}\). Dann gilt
    \(\pi^1_X \neq \pi^2_X\) und \(\pi^1_Y \neq \pi^2_Y\) (wäre eine der
    beiden Gleichheiten erfüllt, ließe sich \(\lambda_3\) trivial wählen
    --- siehe Diskussion in \ref{subsec:intra-scc-invariants-discussion}).
    Das Tripel \((\zeta, \pi^1_{X\cup Y}, \pi^2_{X\cup Y})\) wird in der
    Enumeration in Schritten~\ref{step:enum-z}--\ref{step:enum-xy}
    erreicht, alle drei Realisierbarkeitsabfragen liefern die im Aufruf
    erwarteten Ergebnisse, und der Algorithmus terminiert mit dem
    Gegenbeispiel.
\end{proof}

\subsection{Notwendigkeit der drei Bedingungen}
\label{subsec:intra-scc-invariants-discussion}

Die drei Eigenschaften des Tripels \((\zeta, \pi^1_{X\cup Y},
\pi^2_{X\cup Y})\) --- Z-Gleichheit, X-Unterschied, Y-Unterschied,
gemeinsame Realisierbarkeit und Unerfüllbarkeit der Brücke --- sind alle
notwendig. Wir zeigen an je einem Beispiel, was passiert, wenn eine
davon weggelassen würde.

\paragraph{Z-Gleichheit.}
Die Definition von \(X \bot_{\mathrm{st}} Y \mid Z\) vergleicht nur
stable-Labellings, die auf \(Z\) übereinstimmen. Ließe man Paare mit
\(\lambda_1|_Z \neq \lambda_2|_Z\) als Kandidaten zu, würde die
Brücken-Anfrage in Schritt~\ref{step:bridge} eine Z-Belegung
\(\pi^2_Z\) verlangen, die nicht zu \(\lambda_1|_Z\) passt --- das Paar
testet damit eine andere Frage als die definierte.

\emph{Beispiel.} Sei \(F\) der gerade Zykel
\(a_0\to a_1\to a_2\to a_3\to a_0\) mit den beiden stable-Labellings
\(\lambda^*\) und \(\overline{\lambda^*}\). Mit \(Z=\{a_0\}\),
\(X=\{a_1\}\), \(Y=\{a_2\}\) sind die beiden Labellings auf \(Z\)
verschieden; ohne die Z-Gleichheit als Vorbedingung würde das Paar
\((\lambda^*, \overline{\lambda^*})\) als Kandidat eingespeist, obwohl
es im Sinne der Definition gar nicht vergleichbar ist.

\paragraph{X-Unterschied \texorpdfstring{(\(\pi^1_X \neq \pi^2_X\))}{}.}
Stimmen \(\lambda_1\) und \(\lambda_2\) auf \(X\) überein, so erfüllt
\(\lambda_3 := \lambda_2\) automatisch die Brücken-Bedingung:
\(\lambda_3|_X = \lambda_2|_X = \lambda_1|_X\) und \(\lambda_3|_{Y\cup
Z} = \lambda_2|_{Y\cup Z}\) trivial.
Schritt~\ref{step:bridge} liefert in diesem Fall garantiert SAT.

\emph{Beispiel.} Sei \(F\) der gerade Zykel
\(a_0\to a_1\to a_2\to a_3\to a_0\), \(X=\emptyset\),
\(Y=\{a_0, a_2\}\), \(Z=\emptyset\). Die Stable-Labellings \(\lambda^*\)
und \(\overline{\lambda^*}\) unterscheiden sich auf \(Y\), stimmen auf
\(X\) trivial überein. Schritt~\ref{step:bridge} prüft, ob
\(\overline{\lambda^*}|_Y\) zu einem stable-Labelling erweiterbar
ist --- ja, durch \(\overline{\lambda^*}\) selbst. Kein Gegenbeispiel.

\paragraph{Y-Unterschied \texorpdfstring{(\(\pi^1_Y \neq \pi^2_Y\))}{}.}
Symmetrisch: stimmen \(\lambda_1, \lambda_2\) auf \(Y\) überein, ist
\(\lambda_3 := \lambda_1\) eine Brücke.

\emph{Beispiel.} \(F=(\{a\}, \emptyset)\), einziges stable-Labelling
\(\lambda(a) = \mathit{in}\). Mit \(Y=\{a\}\), \(X=Z=\emptyset\) gibt
es nur ein Paar \((\lambda, \lambda)\), und dieses ist trivialerweise
durch sich selbst zu brücken. Kein Gegenbeispiel.

\paragraph{Realisierbarkeit von \texorpdfstring{\(\lambda_1\)}{lambda_1} und \texorpdfstring{\(\lambda_2\)}{lambda_2}.}
\(\lambda_1, \lambda_2\) müssen tatsächlich existieren. Würde man die
Realisierbarkeitsabfragen in den Schritten~\ref{step:lambda1}
und~\ref{step:lambda2} weglassen, könnte die Brücken-Anfrage UNSAT
liefern, obwohl \(\pi^1_{X\cup Y}\) gar kein gültiges stable-Labelling
fortsetzt --- der Algorithmus würde dann ein ``Pseudo-Gegenbeispiel''
melden, das die Definition nicht erfüllt.

\emph{Beispiel.} \(F=(\{a_0, a_1, a_2\}, \{(a_0, a_1), (a_1, a_2),
(a_2, a_0)\})\) ein ungerader 3-Zykel mit
\(\varLambda_{\mathrm{st}}(F) = \emptyset\). Mit \(X=\{a_0\}\),
\(Y=\{a_1\}\), \(Z=\emptyset\) und \(\pi^1_X(a_0) = \mathit{in},
\pi^1_Y(a_1) = \mathit{in}\) ist die Brücken-Anfrage UNSAT (es gibt
kein einziges stable-Labelling), aber dasselbe gilt für \(\lambda_1\):
das ``Gegenbeispiel'' wäre nur Artefakt des Vakuum-Falls. CI gilt hier
vakuum.

\section{Heuristische Beschleuniger}
\label{sec:intra-scc-heur}

Der Algorithmus aus Abschnitt~\ref{subsec:intra-scc-algo} ist im
Worst Case so teuer wie \(\textit{Ind}_{\mathrm{st}}\) selbst. In der
Praxis lassen sich aber viele Tripel
\((\zeta, \pi^1_{X\cup Y}, \pi^2_{X\cup Y})\) ohne SAT-Aufruf
verwerfen. Im Folgenden notieren wir drei strukturelle Filter, die
sich direkt vor den Realisierbarkeitsabfragen anwenden lassen.

\subsection{Bessere Startpunkte für \texorpdfstring{\(\zeta\)}{zeta}}
\label{subsec:heur-z-pick}

Die Wahl der Reihenfolge in der \(\zeta\)-Enumeration in
Schritt~\ref{step:enum-z} ist für die Korrektheit irrelevant, aber für
das Auffinden eines Gegenbeispiels nicht. Drei Heuristiken bieten sich
an:
\begin{itemize}
    \item \emph{Wenig erzwingende \(\zeta\) zuerst.}
        Berechne für jedes \(\zeta\) den deterministischen Abschluss
        unter \(\mathrm{st}_F(I)\). Belegungen, die viele Argumente
        außerhalb \(Z\) fixieren, lassen wenig Freiheit für
        \(X\)- bzw. \(Y\)-Variationen. Ein \(\zeta\), dessen Abschluss
        \(X\cup Y\) nicht trifft, ist ein vielversprechender Kandidat
        für ein Gegenbeispiel.
    \item \emph{Z-Belegungen, die \(X\cup Y\) komplett fixieren,
        überspringen.}
        Wird durch Propagation jedes Argument in \(X\) oder jedes in
        \(Y\) deterministisch festgelegt, so existiert für dieses
        \(\zeta\) entweder kein \(\pi^1, \pi^2\) mit unterschiedlichen
        \(X\)-Werten, oder keines mit unterschiedlichen \(Y\)-Werten ---
        in beiden Fällen kein Gegenbeispiel. Schritt~\ref{step:enum-xy}
        kann übersprungen werden.
    \item \emph{Vakuum-Fall global prüfen.}
        Ist \(\mathrm{realisierbar}(\emptyset)\) bereits UNSAT, so ist
        \(\varLambda_{\mathrm{st}}(F)=\emptyset\) und
        \(X \bot_{\mathrm{st}} Y \mid Z\) gilt vakuum. Dies kann vor
        der gesamten \(\zeta\)-Enumeration einmalig geprüft werden.
\end{itemize}

\subsection{Filter auf strukturell unmögliche Kombinationen}
\label{subsec:heur-xy-incompat}

In Schritt~\ref{step:enum-xy} werden alle
\((\pi^1_{X\cup Y}, \pi^2_{X\cup Y})\)-Paare aufgezählt. Viele lassen
sich aus rein graphentheoretischen Gründen verwerfen, ohne den
SAT-Solver zu bemühen.

\begin{itemize}
    \item \emph{Direkte Angriffe zwischen \(X\) und \(Y\).} Gilt
        \((x,y)\in R\) mit \(x\in X, y\in Y\), so erzwingt
        \(\pi(x)=\mathit{in}\) sofort \(\pi(y)=\mathit{out}\) in jeder
        stable-Belegung; umgekehrt erlaubt \(\pi(x)=\mathit{out}\)
        beide Werte für \(y\) nur, wenn \(y\) andere Angreifer hat. Pro
        Paar \((x,y)\) lassen sich solche Implikationen vorberechnen
        und auf jede Belegung \(\pi\) direkt prüfen.
    \item \emph{Konflikte innerhalb \(X\cup Y\).} Existieren Angriffe
        \((a,b)\) mit \(a, b\in X\cup Y\) und \(\pi(a)=\pi(b)=
        \mathit{in}\), ist \(\pi\) sofort nicht realisierbar
        (Konfliktfreiheit verletzt).
    \item \emph{Eingebettete gerade Zykel.} Existiert ein induzierter
        gerader Zykel \(D \subseteq A\setminus(X\cup Z)\) ohne externe
        Angreifer, der \(Y\) schneidet, so liefert er deterministisch
        zwei Y-verschiedene Realisierungen für \(\pi^1, \pi^2\) --- die
        beiden alternierenden Färbungen von \(D\). Statt alle
        \(2^{|Y|}\) Belegungen durchzuprobieren genügt es, dieses Paar
        explizit zu konstruieren und nur die Brücken-Anfrage in
        Schritt~\ref{step:bridge} laufen zu lassen.
\end{itemize}

\subsection{Realisierbarkeitsabfrage nutzt nur SAT/UNSAT}
\label{subsec:heur-no-completion}

Die Hilfsroutine \(\mathrm{realisierbar}(\pi)\) ist als reine
Wahrheitswertanfrage formuliert: das SAT-Modell außerhalb
\(X\cup Y\cup Z\) wird vom Algorithmus nicht weiterverwendet. Eine
Implementierung mit inkrementellem SAT-Solver (Assumptions statt
Klausel-Reset) liefert die Antwort in der Größenordnung einer
Konflikt\-analyse pro Aufruf und hält die Solverzustände zwischen den
drei Anfragen warm. Damit fallen die meisten Realisierbarkeitsabfragen
deutlich billiger aus als ein voller QBF-Aufruf, weil weder
Quantorenstruktur noch innere Schleifen wie beim QBF-Solver
durchlaufen werden müssen.

Eine Implementierung kann das Modell aus den Schritten~\ref{step:lambda1}
und~\ref{step:lambda2} zwischenspeichern, falls am Ende ein konkretes
Paar \((\lambda_1, \lambda_2)\) als Zeugnis ausgegeben werden soll.
Für die Entscheidungsfrage selbst sind die Modelle entbehrlich.

